% !TeX root = ../main.tex

\newlecture[][Oct 1, 2026]

\begin{theorem}[Jensen's Inequality][4.5]
    If $f$ is a convex function and $X$ is a random variable such that $f(X)$ is integrable,
    \[f(EX) \leq Ef(X)\]
    The inequality is \dbb{flipped} if $f$ is \dbb{concave}.
    \begin{proof}
    Since $f$ is convex, there exists a function $g$ s.t.
    \begin{align*}
        g(x) &= ax + b \text{ and} & &\db{\Leftarrow \quad g(\cdot) \text{ is a straight line}}\\
        g(x) &\leq f(x) \text{ for all }x \text{ and} & &\db{\Leftarrow \quad g(\cdot) \text{ always below }f(\cdot)}\\
        g(EX) &= f(EX) & &\db{\Leftarrow \quad g(\cdot)=f(\cdot) \text{ at } x=EX}
    \end{align*}
    \begin{center}
    \begin{tikzpicture}[scale=1, every node/.style={font=\small}]
        \draw[->] (-0.2,0) -- (5,0) node[right]{$x$};
        \draw[->] (0,-0.2) -- (0,3.6);
        \draw[thick, domain=0.5:4.3, smooth, variable=\x] plot ({\x}, {0.35*(\x-1.8)^2+1});
        \draw[thick, blue!55!black] (0,0.2825) -- (4.7,1.9275);
        \fill (2.3,1.0875) circle (1.5pt);
        \draw[dotted] (2.3,0) -- (2.3,1.0875);
        \draw[dotted] (0,1.0875) -- (2.3,1.0875);
        \node[below] at (2.3,0) {$E(X)$};
        \node[left, blue!55!black] at (0,1.0875) {$g(EX) = f(EX)$};
        \node[anchor=west] at (4.3,3.2) {$f(\cdot)$};
        \node[anchor=west, blue!55!black] at (4.75,1.93) {\aside{$g(\cdot)$: straight, always below $f(\cdot)$}};
    \end{tikzpicture}
    \end{center}
    \[f(x) \geq g(x) \text{ for all }x\]
    \begin{align*}
        E[f(X)] \geq E[g(X)] &= E[aX+b]\\
        &= aE(X)+b\\
        &= g(E(X))\\
        &= f(E(X))
    \end{align*}
    \end{proof}
\end{theorem}

\begin{example}[Applying the Jensen's Inequality][4.5.1]
    \[[E(X)]^2 \leq E[X^2] \quad f(x) = x^2 \Rightarrow \text{convex}\]
    \[E[\log X] \leq \log [EX] \quad f(x) = \log x \Rightarrow \text{concave}\]
\end{example}

\begin{lemma}[Young's Inequality][4.6]
    For any $p, q > 1$ and $\frac{1}{p}+ \frac{1}{q} = 1$,
    \[xy \leq \frac{x^p}{p} + \frac{y^q}{q} \quad \text{for any }x,y \geq 0 \]
    \begin{proof}
    Let $h_y(x) = \frac{1}{p}x^p + \frac{1}{q}y^q -xy$
    \\To prove $h_y(x) \geq 0$, we take derivatives
    \[h_y'(x) = \frac{dh_y(x)}{dx} = x^{p-1} - y\]
    \[h_y''(x) = \frac{d^2h_y(x)}{dx^2} = (p-1)x^{p-2} \geq 0, \text{ because } x \geq 0\]
    Therefore $h_y'(x) = 0 \Rightarrow x^* = y^{\frac{1}{p-1}}$ is the minimizer.
    \\Plugging in,
    \begin{align*}
        h_y(x^*) &= \frac{1}{p}\left(y^{\frac{1}{p-1}}\right)^p + \frac{1}{q}y^q - y^{\frac{1}{p-1}} \cdot y\\
        &= \frac{1}{p}y^{\frac{p}{p-1}} + \frac{1}{q}y^q - y^{\frac{p}{p-1}}
    \end{align*}
    Since
    \[\frac{1}{p}+\frac{1}{q} = 1 \;\Rightarrow\; \frac{1}{q} = 1 - \frac{1}{p} \;\Rightarrow\; q = \frac{p}{p-1}\]
    substituting $q = \frac{p}{p-1}$ and $\frac{1}{q} = 1 - \frac{1}{p}$:
    \begin{align*}
        h_y(x^*) &= \frac{1}{p}y^{\frac{p}{p-1}} + \left(1-\frac{1}{p}\right)y^{\frac{p}{p-1}} - y^{\frac{p}{p-1}}\\
        &= 0
    \end{align*}
    \[\Rightarrow h_y(x) \geq h_y(x^*) = 0\]
    Therefore $\frac{1}{p}x^p + \frac{1}{q}y^q \geq xy$.
    \\(If $p=q=2$: $\frac{1}{2}x^2 + \frac{1}{2}y^2 \geq xy \;\Leftrightarrow\; x^2+y^2 \geq 2xy$.)
    \end{proof}
\end{lemma}

\begin{theorem}[H\"older's Inequality][4.6]
    If $p,q > 1$ s.t. $\frac{1}{p} + \frac{1}{q} = 1$,
    \[E[|XY|] \leq [E(|X|^p)]^{\frac{1}{p}} \cdot [E(|Y|^q)]^{\frac{1}{q}}\]

    \begin{proof}
    Let
    \[X^* = \frac{|X|}{[E(|X|^p)]^{\frac{1}{p}}} \quad \text{and} \quad Y^* = \frac{|Y|}{[E(|Y|^q)]^{\frac{1}{q}}}\]
    Observe that
    \[E[(X^*)^p] = E[(Y^*)^q] = 1\]
    since
    \[E[(X^*)^p] = E\left[\left(\frac{|X|}{[E(|X|^p)]^{\frac{1}{p}}}\right)^p\right] = E\left[\frac{|X|^p}{E(|X|^p)}\right] = \frac{E(|X|^p)}{E(|X|^p)} = 1\]
    and the same for $Y^*$ with $q$.

    Applying the Young's Inequality
    \[X^*Y^* \leq \frac{(X^*)^p}{p} + \frac{(Y^*)^q}{q}\]
    Take expectations on both sides
    \[E[X^*Y^*]\leq \frac{E[(X^*)^p]}{p} + \frac{E[(Y^*)^q]}{q}\]
    \[E\left[\frac{|X||Y|}{E(|X|^p)^{\frac{1}{p}}E(|Y|^q)^{\frac{1}{q}}}\right] \leq \frac{1}{p}+ \frac{1}{q} = 1\]
    Therefore, re-arrange
    \[E[|X||Y|] \leq E(|X|^p)^{\frac{1}{p}} E(|Y|^q)^{\frac{1}{q}}\]
    \end{proof}
\end{theorem}

\begin{proposition}[Cauchy-Schwarz Inequality][4.6.1]
    \[E[|XY|] \leq \sqrt{E(X^2)E(Y^2)}\]
    \begin{proof}
    Exercise. \aside{(Take $p = q = 2$ in H\"older's inequality: $E[|XY|] \leq [E(X^2)]^{\frac12}[E(Y^2)]^{\frac12}$.)}
    \end{proof}
\end{proposition}

\begin{theorem}[Minkowski's Inequality][4.7]
    If $p>1$, then
    \[[E(|X+Y|^p)]^{\frac{1}{p}} \leq [E(|X|^p)]^{\frac{1}{p}} + [E(|Y|^p)]^{\frac{1}{p}}\]
    \begin{center}
        (The triangle inequality in probability)
    \end{center}
\end{theorem}
\aside{No proof was given in lecture.}

{\large\rb{EXAM UPDATE: CAN BRING CHEATSHEET DOUBLE SIDE HANDWRITTEN!!!!}\par}
